%=================================================================
% HOW TO USE THIS FILE
%=================================================================
% This file contains the manuscript CONTENT only. It is NOT a
% compilable MDPI submission on its own: MDPI requires their own
% class file (mdpi.cls) and a "Definitions/" folder (logos,
% bibliography styles, journal-specific settings), which are NOT
% reproduced here (they are MDPI's own files, not something to
% copy from a third party -- download them fresh, official source
% only, right before finalizing the submission, since MDPI updates
% the template periodically).
%
% Steps to assemble the real submission:
%   1. Go to https://www.mdpi.com/authors/latex and download the
%      current template ZIP (choose the citation style your target
%      journal uses -- ACS/numbered is the default for most
%      MDPI science journals; APA/author-date for a shorter list,
%      Climate uses NUMBERED citations (\cite), not author-date -- see
%      note near \bibliography at the end of this file.
%   2. Unzip it; you will get template.tex + a Definitions/ folder.
%   3. Replace the placeholder journal name below
%      (\documentclass[JOURNALNAME,article,submit,moreauthors]{Definitions/mdpi})
%      with your target journal's exact string, copied from the
%      comment block in MDPI's own template.tex (case-sensitive).
%   4. Copy everything from \Title{...} down to \end{document} in
%      THIS file into MDPI's template.tex, replacing their
%      placeholder content.
%   5. Copy references.bib into the same project folder.
%=================================================================

\documentclass[climate,article,submit,moreauthors]{Definitions/mdpi}
% <-- Confirm 'climate' matches exactly the string in the current official
%     template.tex comment block before compiling (case-sensitive). Chosen
%     over Entropy and over Atmosphere (which already hosts the companion
%     MF-DFA manuscript, so a different journal was preferred for this paper).

%=================================================================
% MDPI internal commands -- fill in once the target journal and
% submission stage are decided; leave as-is for now.
%=================================================================
\firstpage{1}
\makeatletter
\setcounter{page}{\@firstpage}
\makeatother
\pubvolume{1}
\issuenum{1}
\articlenumber{0}
\pubyear{2026}
\copyrightyear{2026}
\datereceived{}
\daterevised{}
\dateaccepted{}
\datepublished{}

%=================================================================
% TITLE -- placeholder, to be written last per the agreed workflow
%=================================================================
\Title{[TITLE PENDING]}

%=================================================================
% AUTHORS
%=================================================================
\TitleCitation{[Short title for citation -- PENDING]}
\newcommand{\orcidauthorA}{0000-0000-0000-0000} % <-- fill in ORCID if available
\Author{Rafael Magallanes-Quintanar $^{1,}$*}
\AuthorNames{Rafael Magallanes-Quintanar}
\AuthorCitation{Magallanes-Quintanar, R.}

\address{%
$^{1}$ \quad Unidad Acad\'emica de Ingenier\'ia El\'ectrica, Programa de Ingenier\'ia en Computaci\'on, Universidad Aut\'onoma de Zacatecas, Zacatecas, Mexico
}

\corres{Correspondence: [email] }

%=================================================================
% ABSTRACT AND KEYWORDS -- placeholders, written last
%=================================================================
\abstract{[ABSTRACT PENDING -- to be written after Discussion is finalized.]}

\keyword{recurrence quantification analysis; cross recurrence plot; drought;
Standardized Precipitation Index; nonlinear time series analysis; synchronization;
Zacatecas; Mexico}

%=================================================================
\begin{document}

% ================================================================
\section{Introduction}
\label{sec:introduction}
[PENDING -- section to be drafted after Discussion.]

% ================================================================
\section{Materials and Methods}
\label{sec:methods}

\subsection{Study area and data}
\label{sec:methods-data}

The study area comprises the state of Zacatecas, north-central Mexico,
partitioned into four physiographic regions (Semi-arid, Highlands,
Mountains, Canyons) via hierarchical cluster analysis of station-level
precipitation series (Section~\ref{sec:methods-regionalization}). This
clustering-based regionalization is not an official administrative or
hydrological classification; rather, it reflects a regionalization approach
that has consistently aligned with the geographic and pluviometric regime of
Zacatecas across the authors' body of work on regional SPI analysis and
forecasting \cite{magallanes2019,magallanes2022,magallanes2023,magallanes2024a,magallanes2024b,magallanes2026forecasting,magallanes2026agrociencia}.
Figure~\ref{fig:stationsmap} shows the spatial distribution of the
stations by climatic region, illustrating how the cluster-based
regionalization aligns with the state's pluviometric regime.

\begin{figure}[H]
\centering
\includegraphics[width=0.85\linewidth]{figures/fig1_stations_map.jpg}
\caption{Meteorological stations used in this study, grouped by climatic region
(Semi-arid, $n=8$; Highlands, $n=12$; Mountains, $n=3$; Canyons, $n=6$; $n=29$
total, after excluding two transitional stations -- see
Section~\ref{sec:methods-regionalization}) obtained via cluster analysis. Inset:
location of the state of Zacatecas within Mexico.}
\label{fig:stationsmap}
\end{figure}

\subsection{Standardized Precipitation Index}
\label{sec:spi}

The Standardized Precipitation Index (SPI) \cite{McKee1993} quantifies the
deviation of observed precipitation from the long-term climatological
distribution for a given accumulation period. Monthly precipitation
accumulated over the timescale of interest was fitted to a two-parameter
gamma distribution, whose probability density function is

\begin{equation}
g(x) = \frac{1}{\beta^{\alpha}\Gamma(\alpha)} x^{\alpha - 1} e^{-x/\beta},
\qquad x > 0,
\label{eq:gamma_pdf}
\end{equation}

where $\alpha > 0$ and $\beta > 0$ are the shape and scale parameters,
respectively, and $\Gamma(\alpha)$ is the gamma function. The parameters
were estimated via the maximum-likelihood approximation of McKee et al.
\cite{McKee1993}:

\begin{equation}
\hat{\alpha} = \frac{1}{4A}\left(1 + \sqrt{1 + \frac{4A}{3}}\right),
\qquad \hat{\beta} = \frac{\bar{x}}{\hat{\alpha}}, \qquad
A = \ln(\bar{x}) - \frac{\sum \ln(x)}{n},
\label{eq:gamma_params}
\end{equation}

where $\bar{x}$ is the sample mean of the precipitation accumulations and
$n$ is the number of observations. Given $\hat{\alpha}$ and $\hat{\beta}$,
the cumulative probability of an observed precipitation amount $x$ is
$G(x) = \int_0^{x} g(x)\,dx$. Because monthly precipitation accumulations
can include zero values with probability $q = m/n$ (where $m$ is the number
of zero-precipitation observations), the actual cumulative probability of
non-exceedance is $H(x) = q + (1 - q)\,G(x)$, which is transformed into the
standard normal variate

\begin{equation}
\mathrm{SPI} = \Phi^{-1}\!\left(H(x)\right),
\label{eq:spi_final}
\end{equation}

where $\Phi^{-1}(\cdot)$ is the inverse standard normal cumulative
distribution function (the probit function). Negative SPI values indicate
drier-than-average conditions and positive values indicate wetter-than-average
conditions. This study uses SPI-12 (the 12-month accumulation, associated
with drought conditions relevant to aquifer recharge and groundwater
levels), computed in \texttt{Python} using a custom implementation of
Equations~\eqref{eq:gamma_pdf}--\eqref{eq:spi_final}, consistent with the
procedure implemented in the widely used \texttt{SPEI} R package
\cite{BegueriaVicenteSerrano2014}.

\subsection{Regionalization}
\label{sec:methods-regionalization}

The 12-month Standardized Precipitation Index (SPI-12; see
Section~\ref{sec:spi})
was computed for each of the 31 stations from the cleaned and imputed
precipitation series. Stations were then grouped into climatically
homogeneous regions by hierarchical cluster analysis (Ward's minimum
variance method, Euclidean distance), computed using the
\texttt{scipy.cluster.hierarchy} and \texttt{scipy.spatial.distance} modules
\cite{Virtanen2020}, applied to the 31 standardized SPI-12 series. The
resulting four-cluster solution was evaluated against the independent
regional classification of \cite{magallanes2019} and against
the geographic and climatological expertise of the authors regarding the
study area. Two stations, Chalchihuites and Concepci\'on del Oro, were
excluded from the final regionalization: both occupied geographically
transitional or ambiguous positions (Chalchihuites lies in a transition zone
toward the Sierra Madre Occidental, and Concepci\'on del Oro is centrally
located within the semi-arid plateau) that resulted in cluster assignments
inconsistent with the well-documented climatology of the region, and, in the
case of Concepci\'on del Oro, with \cite{magallanes2019}.
Excluding these two stations preserved the cluster membership of the
remaining 29 stations. The final regionalization comprised 8 stations in the
Semi-arid region, 12 in Highlands, 3 in Mountains, and 6 in Canyons, with
mean annual precipitation increasing monotonically across regions (409, 464,
538, and 732~mm, respectively), consistent with the expected precipitation
gradient from the semi-arid plateau toward the more humid canyon and
mountain zones of the state. The resulting dendrogram was rendered as a
circular (fan-type) layout using the \texttt{dendrofan} Python package
\cite{MagallanesQuintanar2026Dendrofan}, developed by the authors and
inspired by the visual convention commonly used for hierarchical trees in
the \texttt{ape} R package \cite{Paradis2019}; the clustering computation
itself was performed entirely in \texttt{Python}, using \texttt{SciPy}, as
described above.

The regional SPI-12 series were computed as the arithmetic mean, at each
month, of the station-level SPI-12 values belonging to each region. Because
the 12-month accumulation window requires 11 preceding months of data, the
resulting regional SPI-12 series span 709 months, from December 1965 to
December 2024.

Station-level precipitation series were quality-controlled, gap-filled, and used to
compute the regional SPI-12 series following the regionalization procedure detailed
in Section~\ref{sec:methods-regionalization}; full details of station-level
gamma-distribution fitting and imputation are provided in the companion MF-DFA
manuscript and are not repeated here. The four resulting regional SPI-12 series
(709 months, December 1965--December 2024) constitute the sole input to the
analyses described below.

\subsection{Phase-space reconstruction and embedding parameter selection}
\label{sec:methods-embedding}

Recurrence-based methods require reconstructing the system's phase space from a
scalar time series via time-delay embedding \cite{Takens1981}. For a series
$x(t)$, an embedded state vector is defined as

\begin{equation}
\vec{x}_i = \big(x(i),\, x(i+\tau),\, x(i+2\tau),\, \dots,\, x(i+(m-1)\tau)\big)
\label{eq:embedding}
\end{equation}

where $\tau$ is the delay and $m$ is the embedding dimension.

For each of the four regional series, $\tau$ was estimated as the first local
minimum of the average mutual information (AMI) function, computed over lags
1--24 months using a 16-bin two-dimensional histogram estimator. The embedding
dimension $m$ was estimated via the false nearest neighbors (FNN) method
\cite{Kennel1992}, selecting the smallest $m$ at which the percentage of false
neighbors fell below 5\%, with standard thresholds $R_{tol}=15$, $A_{tol}=2$.
Resulting parameters are reported in Section~\ref{sec:results-embedding}
(Table~\ref{tab:embedding}).

Because the sliding-window synchronization analysis
(Section~\ref{sec:methods-sliding}) operates on much shorter sub-series
(12--36 months) than the full 709-month series, applying the same, comparatively
high-dimensional ($m\approx4$--5, $\tau\approx9$--11 months) embedding to short
windows was found to leave zero usable embedded vectors in many windows
($(m-1)\tau$ alone can consume 27--40 months). For that analysis only, a minimal
fixed embedding ($m=2$, $\tau=1$) was used instead --
see Section~\ref{sec:methods-sliding}.

\subsection{Recurrence plots and RQA measures}
\label{sec:methods-rqa}

For each region, a recurrence plot (RP) was constructed from the embedded
trajectory:

\begin{equation}
R_{i,j} = \Theta(\varepsilon - \lVert \vec{x}_i - \vec{x}_j \rVert)
\label{eq:rp}
\end{equation}

where $\Theta$ is the Heaviside function, $\lVert \cdot \rVert$ is the Euclidean
norm, and $\varepsilon$ is a distance threshold. Rather than fixing $\varepsilon$
to an arbitrary value (e.g., a fixed multiple of the signal's standard
deviation), $\varepsilon$ was selected independently for each region as the
empirical percentile of the pairwise-distance distribution corresponding to a
\textbf{target recurrence rate of RR = 5\%}, following standard recommendations
\cite{Marwan2007} to keep RQA measures comparable across series with different
amplitude scales.

\textbf{Theiler window correction.} Points adjacent in time are trivially close
in phase space whenever the underlying signal is smooth -- which SPI-12, as a
12-month moving average, inherently is. Including such trivial,
autocorrelation-driven matches inflates diagonal-line-based measures (see below)
without reflecting genuine recurrence to a distinct past or future state
\cite{Theiler1986}. To correct for this, a band of half-width $W$ around the
main diagonal ($|i-j| < W$) was excluded from both the $\varepsilon$-calibration
step and the final recurrence matrix. We set $W = \tau$ per region, using the
decorrelation timescale already estimated in
Section~\ref{sec:methods-embedding} as a principled, data-driven choice rather
than an arbitrary one. All RQA measures reported in this manuscript use the
Theiler-corrected recurrence matrices.

From the corrected recurrence matrix, the following standard RQA measures
\cite{Marwan2007} were computed from the distributions of diagonal and
vertical line lengths (minimum line length = 2 for both):

\begin{itemize}
  \item \textbf{RR} (recurrence rate): density of recurrence points,
  $RR = \frac{1}{N^*}\sum_{i,j} R_{i,j}$, where $N^*$ is the number of $(i,j)$
  pairs outside the Theiler-excluded band.
  \item \textbf{DET} (determinism): fraction of recurrence points forming
  diagonal lines of length $\geq L_{min}$.
  \item \textbf{L}: mean diagonal line length.
  \item \textbf{Lmax}: length of the longest diagonal line.
  \item \textbf{ENTR}: Shannon entropy of the diagonal line-length distribution.
  \item \textbf{LAM} (laminarity): fraction of recurrence points forming
  vertical lines of length $\geq V_{min}$ -- an operational measure of states
  "trapped" in a regime.
  \item \textbf{TT} (trapping time): mean vertical line length -- the typical
  duration of a regime before transitioning.
\end{itemize}

\subsection{Cross-Recurrence and the F-synchronization measure}
\label{sec:methods-crp}

Following the approach of Semeraro et al. \cite{Semeraro2020}, pairwise synchronization between regions was
assessed via the Cross Recurrence Plot (CRP):

\begin{equation}
C_{i,j} = \Theta(\varepsilon - \lVert \vec{x}_i - \vec{y}_j \rVert)
\label{eq:crp}
\end{equation}

where $\vec{x}$ and $\vec{y}$ are the embedded trajectories of two different
regional series, with $\varepsilon$ again calibrated to a target RR = 5\% for
the pair. Synchronization was quantified with the \textbf{F-recurrence measure},
the density of points along the CRP's main diagonal (the line of
synchronization):

\begin{equation}
F = \frac{1}{N}\sum_i C_{i,i}
\label{eq:frecurrence}
\end{equation}

$F=1$ indicates complete synchronization; $F \to 0$ indicates desynchronization.
Unlike the single-region RQA measures in Section~\ref{sec:methods-rqa}, $F$ is
computed only along matched time indices $(i,i)$, not across a range of
offsets, and is therefore not subject to the Theiler-window artifact described
above; no correction was applied to this measure.

\subsection{Sliding-window synchronization around documented drought events}
\label{sec:methods-sliding}

To evaluate whether inter-regional synchronization changes around severe
drought events, $F$ was recomputed in 12-month sliding windows (3-month step)
within 36-month periods immediately before and after each of five
independently documented drought years (1996, 2011, 2012, 2021, 2023; sources
and event-specific context in Table~\ref{tab:droughtevents}). Because these
short windows are incompatible with the full-series embedding parameters
(Section~\ref{sec:methods-embedding}), a fixed minimal embedding ($m=2$,
$\tau=1$) was used for this analysis only. Windows yielding fewer than 5
embedded vectors were skipped (this occurred systematically for the 2023
event, for which the series' end date, December 2024, allows only 12 of the
intended 36 post-event months).

For each region pair and event with $\geq 2$ valid windows on both sides, the
pre/post distributions of $F$ were compared with Levene's test (variance
homogeneity) and Welch's $t$-test (difference of means, unequal variances).
Drought years were selected from independently documented sources (national
CONAGUA/SMN records and, where available, Zacatecas-specific peer-reviewed
literature and direct grower testimony; see Table~\ref{tab:droughtevents}),
not chosen post hoc from the data.

\subsection{Multiple-comparison correction}
\label{sec:methods-mcc}

The sliding-window analysis (Section~\ref{sec:methods-sliding}) produces one
test per region pair per event (up to 6 pairs $\times$ 5 events = 30 tests,
reduced to 24 valid tests after excluding windows with insufficient data). At
$\alpha=0.05$, roughly one false positive is expected by chance alone across
this many tests. Both a conservative Bonferroni correction and the less
conservative Benjamini-Hochberg false discovery rate (FDR) procedure were
applied to the resulting $p$-values. One region pair-event
(Semi-arid--Highlands, 2021) was treated as a pre-specified, confirmatory test
rather than an exploratory one, since independent field evidence (grower
testimony) identified 2021 as a Highlands-concentrated event \textit{prior to}
running the analysis; for this test, the FDR-corrected result (not Bonferroni)
is considered the appropriate significance criterion (see
Section~\ref{sec:results-sliding} and Discussion).

\subsection{Software and reproducibility}
\label{sec:methods-software}

All analyses were implemented in Python (NumPy, SciPy, pandas) in a Google
Colab notebook, using a custom implementation of RQA rather than third-party
RQA packages, for full transparency and auditability of every computational
step (embedding, threshold calibration, Theiler-window exclusion, and
line-length extraction). The complete notebook, including the data-loading,
verification, and all diagnostic checks described in
Section~\ref{sec:results}, is provided as supplementary material and archived
on GitHub (\url{https://github.com/tiquis/RQA_SPI}).

% ================================================================
\section{Results}
\label{sec:results}

Figure~\ref{fig:spi12series} shows the four regional SPI-12 series analyzed
throughout this study.

\begin{figure}[H]
\centering
\includegraphics[width=\linewidth]{figures/fig2_regional_spi12_series.png}
\caption{Regional SPI-12 series for Zacatecas (709 months, December
1965--December 2024).}
\label{fig:spi12series}
\end{figure}

\subsection{Embedding parameter selection}
\label{sec:results-embedding}

For each of the four regional series of SPI-12 (Semi-arid, Highlands,
Mountains, Canyons; 709 months, December 1965--December 2024), the embedding
delay ($\tau$) and dimension ($m$) were estimated as described in
Section~\ref{sec:methods-embedding} (Table~\ref{tab:embedding}).

\begin{table}[H]
\caption{Embedding parameters per region.}
\label{tab:embedding}
\newcolumntype{C}{>{\centering\arraybackslash}X}
\begin{tabularx}{\linewidth}{lCC}
\toprule
\textbf{Region} & \boldmath{$\tau$} & \boldmath{$m$} \\
\midrule
Semi-arid & 9  & 5 \\
Highlands & 11 & 4 \\
Mountains & 10 & 5 \\
Canyons   & 11 & 5 \\
\bottomrule
\end{tabularx}
\end{table}

The obtained $\tau$ values (9--11 months) are consistent with the seasonal
periodicity inherent to SPI-12 as a 12-month moving window. The embedding
dimension resulted between 4 and 5 across the four regions, with Highlands as
the only region with $m=4$ instead of $m=5$.

\subsection{Regional RQA measures}
\label{sec:results-rqa}

With the recurrence threshold fixed to achieve a constant recurrence rate
(RR$\approx$0.05) across the four regions, and after excluding a Theiler
window of width $W=\tau$ (per region) around the main diagonal to remove
trivial autocorrelation-driven recurrence \cite{Theiler1986,Marwan2007},
recurrence plots were constructed for each region (Figure~\ref{fig:rps}) and
standard RQA measures were computed (Table~\ref{tab:rqa}).

\begin{figure}[H]
\centering
\includegraphics[width=\linewidth]{figures/fig3_recurrence_plots.png}
\caption{Theiler-corrected recurrence plots for the four regions, with
embedding parameters ($\tau$, $m$) and Theiler window width ($W$) indicated
per panel.}
\label{fig:rps}
\end{figure}

\begin{table}[H]
\caption{RQA measures per region (Theiler-corrected).}
\label{tab:rqa}
\newcolumntype{C}{>{\centering\arraybackslash}X}
\begin{tabularx}{\linewidth}{lCCCCCCC}
\toprule
\textbf{Region} & \textbf{RR} & \textbf{DET} & \textbf{L} & \textbf{Lmax} & \textbf{ENTR} & \textbf{LAM} & \textbf{TT} \\
\midrule
Semi-arid & 0.050 & 0.376 & 3.346 & 21 & 1.582 & 0.854 & 4.220 \\
Highlands & 0.050 & 0.380 & 3.328 & 21 & 1.582 & 0.864 & 4.458 \\
Mountains & 0.050 & 0.374 & 3.209 & 21 & 1.515 & 0.849 & 4.188 \\
Canyons   & 0.050 & 0.382 & 3.708 & 30 & 1.751 & 0.863 & 4.895 \\
\bottomrule
\end{tabularx}
\end{table}

\textit{Methodological note (Theiler window):} an initial run without
excluding the near-diagonal band produced Lmax values of 143--343 months,
driven in the most extreme case (Mountains) entirely by an offset-1 diagonal
-- i.e., consecutive months trivially resembling each other due to the
smoothness of SPI-12 as a 12-month moving average, not genuine recurrence to
a distinct past state. Following the standard correction of
Theiler \cite{Theiler1986}, a band of width $W=\tau$ was excluded from all RQA
computations; Table~\ref{tab:rqa} reflects the corrected values used
throughout this manuscript.

Canyons presents the highest values of DET, L, Lmax, ENTR, LAM, and TT among
the four regions -- this ranking is unchanged by the Theiler-window
correction, indicating it is a robust feature of the data rather than an
artifact of the near-diagonal band. Differences between regions remain modest
in absolute terms (e.g., LAM between 0.849 and 0.864), consistent with shared
regional climatic forcings across all four regions.

\textit{Verification of Lmax (Mountains):} the longest diagonal line
(length = 21 months) after correction links two temporally distant wet
episodes -- July 1985--March 1987 and March 2002--November 2003, offset by
200 months ($\approx$16.7 years) -- both with positive mean SPI-12 (0.30 and
0.36, respectively) and healthy within-segment variability (std 0.63 and
0.39; 21/21 unique values in both segments, ruling out a data artifact). A
direct Pearson correlation between the two segments (aligned by relative
month) confirms a moderate-to-strong positive relationship ($r=0.53$),
consistent with genuine recurrence between two similar wet regimes rather
than a methodological artifact (Figure~\ref{fig:lmaxdiag}).

\begin{figure}[H]
\centering
\includegraphics[width=0.75\linewidth]{figures/fig5_lmax_diagnostic_mountains.png}
\caption{Mountains region: the two temporally distant segments (Segment A,
1985--1987; Segment B, 2002--2003) forming the longest diagonal line
(Lmax = 21 months) in the Theiler-corrected recurrence plot, aligned by
relative month since segment start.}
\label{fig:lmaxdiag}
\end{figure}

\subsection{Synchronization between regions (Cross-Recurrence)}
\label{sec:results-crp}

The F-synchronization measure (density of the identity line in the Cross
Recurrence Plot; \cite{Semeraro2020}) was computed for the six pairwise
region combinations (Table~\ref{tab:frecurrence}).

\begin{table}[H]
\caption{Pairwise F-recurrence (synchronization) between regions.}
\label{tab:frecurrence}
\newcolumntype{C}{>{\centering\arraybackslash}X}
\begin{tabularx}{0.6\linewidth}{lC}
\toprule
\textbf{Region pair} & \textbf{F} \\
\midrule
Semi-arid -- Highlands & \textbf{0.639} \\
Highlands -- Mountains & 0.573 \\
Semi-arid -- Mountains & 0.461 \\
Mountains -- Canyons   & 0.460 \\
Highlands -- Canyons   & 0.352 \\
Semi-arid -- Canyons   & 0.287 \\
\bottomrule
\end{tabularx}
\end{table}

\begin{figure}[H]
\centering
\includegraphics[width=0.6\linewidth]{figures/fig4_frecurrence_heatmap.png}
\caption{Heatmap of pairwise F-recurrence (synchronization) between regions.
Diagonal cells (self-synchronization, $F=1$ by definition) are marked
``diag''.}
\label{fig:fheatmap}
\end{figure}

Semi-arid and Highlands constitute, by a clear margin, the most synchronized
pair -- consistent with their geographic contiguity in the inter-range basin
between the Sierra Madre Oriental and Occidental. Semi-arid and Canyons
constitute the least synchronized pair, consistent with distinct orographic
precipitation regimes between the two regions. The heatmap in
Figure~\ref{fig:fheatmap} summarizes this pattern, showing a clear gradient
of decreasing synchronization from geographically contiguous to more
distant/orographically distinct region pairs.

\subsection{Synchronization before/after documented drought events}
\label{sec:results-sliding}

The change in the $F$ measure was evaluated before and after five
independently documented drought events (Table~\ref{tab:droughtevents}), for
the six region-pair combinations. Of the resulting 24 tests (2023 was
excluded due to insufficient post-event data within the series range),
correction for multiple comparisons was applied (Table~\ref{tab:sliding}).

\begin{table}[H]
\caption{Region pairs with statistically significant synchronization change around a drought event (after multiple-comparison correction).}
\label{tab:sliding}
\newcolumntype{C}{>{\centering\arraybackslash}X}
\begin{tabularx}{\linewidth}{lCCCCCC}
\toprule
\textbf{Pair} & \textbf{Event} & \textbf{F before} & \textbf{F after} & \boldmath{$p$} \textbf{(raw)} & \boldmath{$p$} \textbf{(Bonf.)} & \textbf{Sig. FDR} \\
\midrule
Semi-arid--Highlands & 2012 & 0.020 & 0.323 & $2.27\times10^{-8}$ & $5.45\times10^{-7}$ & \checkmark \\
Highlands--Canyons   & 2021 & 0.010 & 0.253 & $2.45\times10^{-3}$ & $5.89\times10^{-2}$ & \checkmark \\
Semi-arid--Highlands & 2021 & 0.313 & 0.091 & $8.10\times10^{-3}$ & $1.94\times10^{-1}$ & \checkmark \\
\bottomrule
\end{tabularx}
\end{table}

The three remaining pair-events with uncorrected $p<0.05$
(Semi-arid--Canyons 2012, $p=0.025$; Highlands--Canyons 2011, $p=0.034$;
Mountains--Canyons 2011, $p=0.036$) did not survive any correction and are
reported as inconclusive trends, not as findings. The 1996 event did not
produce significant changes in any region pair ($p>0.05$ across all five
pairs evaluated).

The two events documented as being of systemic/national scale within the
period with sufficient data (2011--2012) are associated with a marked
increase in synchronization in the geographically best-connected pair
(Semi-arid--Highlands). The 2021 event, documented via grower testimony as
regionally concentrated in Highlands municipalities (Sombrerete, Valparaíso,
Río Grande), produced a divergent response across pairs: synchronization in
Highlands--Canyons but desynchronization in Semi-arid--Highlands -- the pair
that under normal conditions is the most strongly coupled
(Section~\ref{sec:results-crp}).

\begin{table}[H]
\caption{Documented drought events used in the sliding-window analysis.}
\label{tab:droughtevents}
\newcolumntype{C}{>{\centering\arraybackslash}X}
\begin{tabularx}{\linewidth}{lXl}
\toprule
\textbf{Event} & \textbf{Documentation} & \textbf{Scale} \\
\midrule
1996 & Worst drought of the 1990s nationally; paralyzed agricultural exports; mass livestock sacrifice \cite{TVAztecaSequias} & Systemic/national \\
2011 & Record affected area since satellite records began; 95.12\% of national territory affected, driven by La Ni\~na \cite{SMNsequia2026} & Systemic/national, continental \\
2011--2012 & Worst drought in 70 years in Zacatecas; $\sim$13-month dry spell; $>$75,000 head of livestock lost \cite{Scielo2016Sequia,Scielo2018Sequia,AliatConexxion2012} & Zacatecas, systemic \\
2021--2022 & $>$85\% of country critical (Apr. 2021); urban water crisis in northern states by 2022 \cite{SMNsequia2026} & Continental/national \\
2021 & Critical ranching year in Sombrerete, Valparaíso, R\'io Grande (Highlands); direct grower testimony & Zacatecas, regionally concentrated \\
2023--2024 & 76\% of national territory D0--D4; reservoirs at historic lows \cite{Informador2024Sequia,WikipediaEscasezHidrica,Ganaderia2024Sequia} & Systemic/national \\
2023 & 50 of 58 Zacatecas municipalities in extreme drought (CONAGUA); not statistically testable, data ends Dec. 2024 & Zacatecas, systemic \\
\bottomrule
\end{tabularx}
\end{table}

% ================================================================
\section{Discussion}
\label{sec:discussion}

\subsection{Regional Persistence and the Precipitation Gradient}
\label{sec:disc-persistence}

Canyons consistently showed the highest values across all diagonal- and
vertical-line RQA measures (DET, L, Lmax, ENTR, LAM, TT; Section
\ref{sec:results-rqa}), and this ranking did not change after the
Theiler-window correction was applied, indicating that it is a genuine
feature of the regional dynamics rather than an artifact of the near-diagonal
band. This finding is broadly consistent with the precipitation gradient
reported in Section~\ref{sec:methods-regionalization}: mean annual
precipitation increases monotonically from Semi-arid (409~mm) to Highlands
(464~mm), Mountains (538~mm), and Canyons (732~mm). A plausible explanation
is that regions with stronger, more topographically anchored orographic
forcing -- such as Canyons -- experience a more regular seasonal
precipitation cycle, which may translate into more laminar, self-similar
SPI-12 trajectories (higher LAM and TT); by contrast, Semi-arid, dominated by
more sporadic convective precipitation over open plains, may be more
susceptible to short, less regular excursions that reduce apparent
persistence. This interpretation remains a hypothesis rather than an
established mechanism, and it should be noted that the absolute differences
between regions were modest (e.g., LAM ranged only from 0.849 to 0.864),
consistent with all four regions sharing common synoptic-scale climatic
forcings despite their differing local topography.

This finding can also be considered alongside two related analyses conducted by the
authors on the same four regions. A multifractal detrended fluctuation analysis (MF-DFA)
of the same regional SPI-12 series, currently unpublished, found that Canyons exhibited
the smallest degree of multifractality among the four regions, while Semi-arid exhibited
the largest -- broadly consistent with the comparatively higher RQA-based persistence
(LAM, TT) found for Canyons here, insofar as both measures point toward more regular,
less heterogeneous short-term dynamics in that region relative to Semi-arid. Separately,
a progression of published forecasting studies by the authors, applying NARX networks
\cite{magallanes2022}, an N-HiTS/LSTM comparison \cite{magallanes2024a}, and an
AutoML approach \cite{magallanes2024b} to the same four regions, found that Canyons
was, in apparent contrast, among the most difficult regions to forecast at short-to-medium
horizons in two of these three studies: the lowest test-period correlation of the four
regions in \cite{magallanes2022} ($R=0.844$, versus 0.953 for Semi-arid) and in
\cite{magallanes2024b} ($R=0.933$, versus 0.964 for Highlands), accompanied in both
studies by the strongest under-prediction bias among the four regions (58.6\% and 62\%
probability of under-prediction, respectively). A related, currently unpublished analysis
comparing Transformer-based architectures across the same four regions similarly found
Canyons to have the weakest predictive-interval calibration of the four regions, though
with the largest, statistically confirmed benefit from incorporating a macroclimatic
covariate. Taken together, these complementary characterizations suggest that the
regularity of Canyons captured by standard persistence and multifractal measures does
not consistently translate into short-horizon predictability, plausibly because its dynamics
are shaped by comparatively rare but large-magnitude regime transitions -- consistent with
the pronounced swings between severe drought and rapid recovery visible in its regional
SPI-12 series (Figure~\ref{fig:spi12series}) -- rather than by the more frequent,
moderate-magnitude fluctuations that standard forecasting metrics are most sensitive to.
This apparent dissociation between long-term regularity and short-term predictability is,
to the authors' knowledge, not yet formally established, and constitutes a natural
direction for future work relating RQA-based persistence measures directly to forecast
error.

\subsection{Geographic Basis of Inter-Regional Synchronization}
\label{sec:disc-sync}

Semi-arid and Highlands were, by a clear margin, the most synchronized pair
(F = 0.639; Section~\ref{sec:results-crp}), whereas Semi-arid and Canyons
were the least synchronized (F = 0.287). This pattern is consistent with the
physical geography of the state: Semi-arid and Highlands occupy a shared
inter-range basin between the Sierra Madre Oriental and the Sierra Madre
Occidental, and are therefore likely to be exposed to similar large-scale
moisture-transport pathways and topographic sheltering, whereas Canyons sits
in a distinct orographic setting more conducive to localized convective
precipitation. The use of a cross-recurrence-based synchronization measure to
detect such spatial coherence follows the general approach of
\cite{Semeraro2020}, who used the analogous F-recurrence measure to assess
synchronization between vegetation-index time series in a different
ecological context. The gradient of decreasing synchronization observed here
-- from geographically contiguous to more topographically distinct region
pairs (Figure~\ref{fig:fheatmap}) -- suggests that RQA-based synchronization
measures may be sensitive to sub-state hydroclimatic structure that is not
captured by regional aggregation alone.

\subsection{Systemic versus Regionally Concentrated Drought Signatures}
\label{sec:disc-events}

The sliding-window analysis (Section~\ref{sec:results-sliding}) identified
three region pair-events for which the change in synchronization survived
correction for multiple comparisons: an increase in Semi-arid--Highlands
synchronization around the 2011--2012 drought, and both an increase in
Highlands--Canyons synchronization and a decrease in Semi-arid--Highlands
synchronization around the 2021 drought. These two episodes were not
expected, a priori, to produce the same signature: 2011--2012 is documented
as a systemic, near-nationwide event, whereas 2021 is documented via direct
grower testimony as concentrated in Highlands municipalities. The observed
contrast -- convergence of an already well-coupled pair under a systemic
event, versus a divergent, pair-dependent response under a regionally
concentrated event -- is consistent with the general expectation that the
spatial reach of a climatic anomaly, and not only its severity, shapes its
synchronization signature. A broadly similar principle was reported by
\cite{Semeraro2020}, who found that a large-scale ENSO event produced a
temporary disruption in vegetation-index synchronization in the Amazon,
whereas the comparably large-scale AMO did not, implying that not all
large-scale climatic anomalies leave the same imprint on local
synchronization structure. It should be emphasized that only 3 of 24 tests
survived correction, and that this pattern is therefore best regarded as
hypothesis-generating rather than conclusive; the 2023 drought, independently
documented as one of the most severe in the study period, could not be
evaluated within the current windowing scheme because the series ends in
December 2024, and represents a natural candidate for confirming or
revising this interpretation once additional post-event data become
available.

\subsection{Methodological Considerations}
\label{sec:disc-methods}

Two methodological findings of this study are likely to generalize beyond
the present case. First, the Theiler-window correction
(Section~\ref{sec:results-rqa}) was not a minor refinement: uncorrected Lmax
values were inflated by up to an order of magnitude, driven by trivial
autocorrelation between adjacent months rather than genuine recurrence. This
is expected whenever RQA is applied to a moving-average index such as
SPI-12, and, consistent with the original recommendation of
\cite{Theiler1986} and its treatment in \cite{Marwan2007}, should be treated
as a required step -- rather than an optional one -- in future applications
of RQA to smoothed hydroclimatic series at any accumulation timescale.
Second, the embedding parameters appropriate for the full 709-month series
proved entirely unsuitable for the 12--36-month windows used in the
sliding-window analysis, since $(m-1)\tau$ alone could consume most or all
of a short window. This scale-dependence of embedding parameters is a
general property of phase-space reconstruction and, in the specific context
of short-window drought-event analysis, implies that a separate, minimal
embedding is generally necessary rather than a simplification adopted for
convenience.

\subsection{Limitations}
\label{sec:disc-limitations}

Several limitations should be considered when interpreting these results.
First, the four-region classification, although grounded in a
cluster-analysis approach that has repeatedly aligned with the geographic
and pluviometric regime of Zacatecas across the authors' prior work
\cite{magallanes2019,magallanes2022,magallanes2023,magallanes2024a,magallanes2024b,magallanes2026forecasting,magallanes2026agrociencia},
is a data-driven regionalization rather than an official administrative or
hydrological one, and alternative regionalizations might reveal different
synchronization structure. Second, a related, currently unpublished analysis
by the authors, comparing deep-learning forecasting configurations across
the same four regions, found that Semi-arid was the only region whose
station network did not meet a strict multi-variable completeness
criterion, requiring a more flexible station-matching scheme due to
instrumentation gaps; this same region also showed the lowest RQA-based
persistence (DET, LAM, TT) in the present study. It is therefore possible
that part of the regional differentiation reported here reflects
heterogeneity in station-network density or completeness across regions,
rather than climatic differences alone, and this potential confound cannot
be fully disentangled with the present data. Third, the sliding-window
synchronization
analysis (Section~\ref{sec:methods-sliding}) used a minimal, fixed embedding
distinct from the one used for the full-series RQA measures; the two
therefore capture related but not strictly comparable aspects of the
system's dynamics, and this distinction should be kept in mind when relating
findings across Sections~\ref{sec:results-rqa} and
\ref{sec:results-sliding}. Fourth, the 2023 drought event could not be
statistically evaluated because the series ends in December 2024, leaving
the most recent, and independently documented as one of the most severe,
episode in the study period outside the scope of formal testing. Finally,
SPI-12 captures meteorological drought as reflected in precipitation
anomalies but does not directly represent soil moisture, agricultural
impact, or water-management outcomes; where available, the qualitative
consistency between the statistical findings and independently documented
regional impacts (Table~\ref{tab:droughtevents}) provides some
triangulation, but a more formal joint analysis with impact-based drought
indicators remains a direction for future work.

% ================================================================
\section{Conclusions}
\label{sec:conclusions}
[PENDING]

% ================================================================
% BACK MATTER -- MDPI requires these sections even if brief.
% Fill in before submission.
% ================================================================

\vspace{6pt}
\authorcontributions{[PENDING -- CRediT taxonomy statement.]}
\funding{[PENDING]}
\institutionalreview{Not applicable.}
\informedconsent{Not applicable.}
\dataavailability{The raw precipitation data used to compute the regional SPI-12
series are part of the ``Proyecto de bases de datos climatológicos,'' maintained by
the Comisión Nacional del Agua (CONAGUA), the official Mexican institution
responsible for climatic and meteorological data record-keeping, and are publicly
available at
\url{https://smn.conagua.gob.mx/es/climatologia/informacion-climatologica/informacion-estadistica-climatologica}.
The complete analysis code, including the data-loading, verification, and all
diagnostic checks described in Section~\ref{sec:results}, is openly available at
\url{https://github.com/tiquis/RQA_SPI}.}
\acknowledgments{[PENDING]}
\conflictsofinterest{The authors declare no conflict of interest.}

%=================================================================
% References -- Climate uses NUMBERED citation style (\cite), already
% applied throughout this file. Re-verify against the current
% instructions-for-authors page before submission, in case MDPI has
% changed this journal's convention.
%=================================================================
\reftitle{References}
\bibliography{references}

\end{document}
